\section{PaLM、Flan-PaLM 与 Med-PaLM}

评价框架就位后，课程开始比较三层模型：PaLM 提供大规模基础能力，Flan-PaLM 通过多任务 instruction tuning 改善指令遵循，Med-PaLM 再用医学 exemplar 与 instruction prompt tuning 做领域对齐。每一步改变的对象不同，理解这条链比记住一个排行榜数字更重要。

\subsection{规模与 instruction tuning 的分工}

PaLM 页展示 Pathways Language Model 的规模与架构背景；结果页说明扩大模型在多种任务上带来明显增益。Flan-PaLM 页则强调 instruction fine-tuning：模型在大量用自然语言描述的任务上训练，学习“读懂任务并按要求作答”，而不只是继续预测下一 token。

\lecturefigure{slide-23-palm-architecture-scale.jpg}{PaLM：大规模基础语言模型}{官方录像恢复幻灯片 V023；Stanford CS25 Lecture 17}

\lecturefigure{slide-24-palm-benchmark-results.jpg}{PaLM 在通用 benchmark 上的规模收益}{官方录像恢复幻灯片 V024；Stanford CS25 Lecture 17}

\lecturefigure{slide-25-flan-palm.jpg}{Flan-PaLM：通过 instruction tuning 学习任务遵循}{官方录像恢复幻灯片 V025；Stanford CS25 Lecture 17}

\begin{knowledgebox}{术语消化：三种“调优”}
\term{continued pretraining} 继续优化模型参数以吸收领域语料；\term{instruction tuning} 用许多显式任务与示例更新模型参数，使其遵循自然语言指令；\term{instruction prompt tuning} 冻结大模型，只学习一小段连续 soft prompt，并配合医学 exemplar 引导模型。三者的成本、可控性和证据都不同。
\end{knowledgebox}

\subsection{医学选择题结果与 Scaling 曲线}

三张 benchmark 页分别覆盖 MedQA、MedMCQA/PubMedQA 与 MMLU clinical topics。应先看相对趋势：更大的 Flan-PaLM 显著超过早期专用基线，表明通用模型的规模与 instruction tuning 能迁移到医学选择题；再看限制：题目通常封闭、答案短，而且准确率不检查模型给出的解释是否安全。

\lecturefigure{slide-26-medqa-sota.jpg}{MedQA（USMLE 风格）上的 state-of-the-art 结果}{官方录像恢复幻灯片 V026；Stanford CS25 Lecture 17}

\lecturefigure{slide-27-medmcqa-pubmedqa-results.jpg}{MedMCQA 与 PubMedQA 结果比较}{官方录像恢复幻灯片 V027；Stanford CS25 Lecture 17}

\lecturefigure{slide-28-mmlu-clinical-results.jpg}{MMLU 临床主题上的模型结果}{官方录像恢复幻灯片 V028；Stanford CS25 Lecture 17}

Scaling plots 页把多个模型规模画在横轴上，纵轴是不同医学任务的准确率。主要趋势是规模增大时性能上升，但不同任务的斜率和饱和点不同；这说明 scale 是强变量，却不是唯一变量。若两个模型拥有不同 instruction-tuning 数据，曲线差异也不能全部归因于参数量。

\lecturefigure{slide-29-scaling-plots.jpg}{医学问答中的 Scaling 趋势}{官方录像恢复幻灯片 V029；Stanford CS25 Lecture 17}

比较 scaling 曲线时还要控制训练与推理预算。更大的模型可能看过更多 token、使用更强 instruction mixture，也可能在 few-shot prompt 中获得更多示例；如果这些条件不同，横轴“参数量”并不是唯一变化。医疗场景还会增加延迟、成本与隐私约束：一个云端超大模型在 benchmark 上多几个百分点，未必胜过可在医院受控环境运行、能够引用院内指南并稳定拒答的小模型。课堂提出“还需要 clinical LM 吗”正是提醒我们，Pareto frontier 至少包含准确率、长答案风险、推理成本、数据治理和可更新性，而非单一排行榜。

\begin{warningbox}{读图时不能推出什么}
一条向上的 scaling curve 不能证明模型形成了可靠临床推理，也不能证明继续增大参数必然消除伤害与偏见。它只说明在当前数据、提示和评分规则下，规模与任务准确率相关；长答案证据必须另行测量。
\end{warningbox}

\subsection{Selective prediction：允许模型拒答}

Selective prediction（选择性预测）让模型只回答高置信度样本。设样本集合为 $D$，模型给出答案 $\hat y_i$ 与置信分数 $s_i$，阈值为 $\tau$。覆盖率与条件准确率分别为
$$
\operatorname{coverage}(\tau)=\frac{1}{|D|}\sum_{i\in D}\mathbf{1}[s_i\ge\tau],
\qquad
\operatorname{acc}(\tau)=
\frac{\sum_i\mathbf{1}[s_i\ge\tau]\mathbf{1}[\hat y_i=y_i]}
{\sum_i\mathbf{1}[s_i\ge\tau]}.
$$
其中 $D$ 表示评测样本，$s_i$ 表示第 $i$ 个回答的置信度，$\tau$ 表示拒答阈值，$\mathbf{1}[\cdot]$ 是指示函数，$y_i$ 是正确答案。提高 $\tau$ 通常降低 coverage、提高保留下来的准确率，但前提是置信度与真实错误相关。

\lecturefigure{slide-30-selective-prediction.jpg}{Selective prediction：覆盖率与准确率的权衡}{官方录像恢复幻灯片 V030；Stanford CS25 Lecture 17}

\begin{lstlisting}[caption={选择性预测的阈值化流程（示意）}]
def selective_answer(question, threshold):
    answer, confidence = model.generate_with_confidence(question)
    if confidence < threshold:
        return "ABSTAIN"
    return answer
\end{lstlisting}

\begin{importantbox}{拒答能力的前提}
abstention 只有在置信度校准有效、拒答能够交给更可靠的人或系统、且低覆盖率可被业务接受时才有价值。医学系统还必须检查“模型最自信的错误”，因为错误校准会把风险样本保留下来。
\end{importantbox}

选择性预测还需要定义拒答后的动作。如果低置信度答案只是返回空白，用户可能转向更不可靠的信息源；如果系统自动升级给 clinician，就要测量升级队列、响应时间和工作负担。阈值也不能只在静态 benchmark 上选择，因为真实问题分布会漂移，新的药物、指南和罕见病例可能系统性地产生过度自信。一个完整监控回路应记录 coverage、错误率、校准误差、拒答原因和人工接管结果，并按科室或人群分层检查。这样才能知道“少答一些”是否真的把风险转移给更可靠的流程，而不是把困难用户排除在服务之外。

\subsection{为什么仍需要医学域对齐}

Human evaluation limitations 页承认开放答案难以规模化评估；下一页则回答为何不能停在 Flan-PaLM：通用模型虽然选择题很强，长答案仍可能偏离医学共识、遗漏关键信息或不适合患者阅读。医学数据与专家 rubric 不是为了重新训练全部世界知识，而是为了改变回答行为。

\lecturefigure{slide-31-human-eval-limitations.jpg}{人类长答案评价的规模与成本限制}{官方录像恢复幻灯片 V031；Stanford CS25 Lecture 17}

\lecturefigure{slide-32-need-domain-data.jpg}{为什么仍需要医学领域数据与对齐信号}{官方录像恢复幻灯片 V032；Stanford CS25 Lecture 17}

\subsection{Instruction prompt tuning 与 Med-PaLM 流程}

Prompt tuning 页显示：大模型参数保持冻结，训练的是一组连续 prompt vectors；医学 exemplar 则为任务提供格式与语义示范。将二者放在一起，得到 Med-PaLM 的对齐流程。它不是课堂旧稿所写的 RLHF，也没有在课上引入双语 fairness prompt 或每周 clinician review 回灌；可信重建应只保留幻灯片和论文支持的步骤。

\lecturefigure{slide-33-prompt-tuning.jpg}{Instruction prompt tuning：冻结模型并学习软提示}{官方录像恢复幻灯片 V033；Stanford CS25 Lecture 17}

\lecturefigure{slide-34-medpalm-pipeline.jpg}{将医学 exemplar 与 prompt tuning 组合为 Med-PaLM}{官方录像恢复幻灯片 V034；Stanford CS25 Lecture 17}

设冻结语言模型为 $f_{\theta}$，可学习软提示为 $P\in\mathbb{R}^{m\times d}$，输入 token embedding 为 $X$，目标答案为 $Y$。训练目标可写为
$$
P^{*}=\arg\min_{P}\; -\sum_{t=1}^{|Y|}\log p_{\theta}(Y_t\mid P,X,Y_{<t}).
$$
其中 $\theta$ 是冻结参数，$m$ 是软提示长度，$d$ 是 embedding 维度，$Y_t$ 是第 $t$ 个目标 token。更新量集中在 $P$，因此参数效率高；但软提示学到什么仍取决于 exemplar、训练任务与评价标准。

\teachervoice{讲者把这一段称为“putting it all together”：基础模型提供广泛知识，instruction tuning 提供任务遵循，医学 prompt tuning 与 exemplar 提供领域行为。课堂没有把任何一步说成安全保证，后面仍要逐轴与 clinician answer 对比。}

\begin{warningbox}{不要把 soft prompt 写成可读的系统提示}
连续 prompt vectors 是模型 embedding 空间里的可学习参数，不一定对应人类可读的词。它与在聊天框中追加“请谨慎回答”的离散提示不同，也不自动提供可解释的医学规则。
\end{warningbox}

从可复现性看，这条 pipeline 至少需要记录四类配置：基础模型 checkpoint、instruction-tuning 版本、医学 exemplar 集合和 soft-prompt 参数。只发布最后一个模型名，无法判断改进来自规模、任务混合还是示例选择。尤其在医学问答中，few-shot exemplar 可能泄漏答案格式或隐含标签分布；因此应在独立问题集上重新选择阈值，并报告 prompt sensitivity。参数高效并不等于评价简单：虽然只更新少量向量，模型行为仍可能在不同措辞、问题长度和用户背景下明显变化。讲义把这些组件拆开，是为了让读者能够定位失败，而不是把所有行为归因于“Med-PaLM”这个单一标签。

\subsection{本章小结}

PaLM、Flan-PaLM 与 Med-PaLM 分别贡献基础规模、通用指令遵循和医学域行为。选择题与 scaling 结果证明大模型具有强迁移能力，selective prediction 提供拒答机制，而 instruction prompt tuning 以较少可训练参数完成领域对齐；这些机制仍必须通过长答案的人类评价检验。

